\documentclass[a4paper]{report}
\usepackage{csquotes}
\usepackage[colorlinks]{hyperref}

\title{Course Evaluation Summary for vetcyb24p2 and vetcyb25p4}
\author{Daniel Bosk and GPT-5}

\begin{document}
\maketitle
\tableofcontents

\chapter{vetcyb24p2}

\section{Quantitative Summary}

\subsection{555894: How much time did you spend on the teaching of the course?}

That is until the teaching ended before the exam period. (The expected workload 
is 80 hours for the whole course, so roughly 8 h/week for 10 weeks. The 
teaching period is 7 weeks.)
Don't include the time for watching the TAMOS lectures, I count that time to 
the TAMOS course.

Total responses: 45\\

\subsubsection{Response distribution:}
\begin{itemize}
  \item 6-8 h/week: 11 (24.4\%)
  \item 4-6 h/week: 10 (22.2\%)
  \item 2-4 h/week: 9 (20.0\%)
  \item 8-10 h/week: 7 (15.6\%)
  \item 0-2 h/week: 6 (13.3\%)
  \item 10-12 h/week: 2 (4.4\%)
\end{itemize}

\subsection{555895: Below you'll find a few statements about the prep material. Check all that you consider true.}

Total responses: 44\\
Total selections: 114\\

\subsubsection{Option distribution:}
\begin{itemize}
  \item It was too many papers to read for each seminar.: 25 (56.8\%)
  \item It was helpful to see others' annotations in the papers that we had to read.: 24 (54.5\%)
  \item It's a good idea to annotate the TAMOS lectures from a security perspective.: 21 (47.7\%)
  \item Annotating the TAMOS lectures would be more useful with more guidance.: 13 (29.5\%)
  \item It would be helpful to not see the others' annotations in the papers until done with one's own annotations.: 11 (25.0\%)
  \item It was a suitable number of papers to read for each seminar.: 11 (25.0\%)
  \item Some topics would have benefited from more papers as example methodologies.: 9 (20.5\%)
\end{itemize}

\subsection{555896: Below you'll find a few statements about the prep material.Check all that you consider true.}

Total responses: 39\\
Total selections: 69\\

\subsubsection{Option distribution:}
\begin{itemize}
  \item Reading the papers was good prep for the seminars.: 33 (84.6\%)
  \item The reflection instructions can be improved to better prepare us for the seminars.: 18 (46.2\%)
  \item Writing the reflections was good prep for the seminars.: 16 (41.0\%)
  \item Some topics should have had more papers to better prepare for the seminar.: 2 (5.1\%)
\end{itemize}

\subsection{555898: Below you'll find a few statements about the seminars. Check all that you consider true.}

Total responses: 40\\
Total selections: 83\\

\subsubsection{Option distribution:}
\begin{itemize}
  \item The seminars helped me deepen my understanding for the science of security.: 27 (67.5\%)
  \item The seminars helped me better understand how the methods learned in TAMOS can be used in security.: 24 (60.0\%)
  \item It was helpful to discuss the topics with classmates during the seminars.: 23 (57.5\%)
  \item Some that I discussed with during a seminar weren't prepared so they couldn't contribute anything to the discussion.: 9 (22.5\%)
\end{itemize}

\subsection{556031: Below you'll find a few statements about the project.Check 
all that you consider true.}

Total responses: 44\\

\subsubsection{Response distribution:}
\begin{itemize}
  \item I didn't have to do the project since I did all prep and participated : 24 (54.5\%)
  \item I missed a seminar so I had to do the project.: 11 (25.0\%)
  \item I never planned on participating during the course\, I was aimin for t: 7 (15.9\%)
  \item I missed a seminar so I had to do the project.,I never planned on part: 2 (4.5\%)
\end{itemize}

\section{Qualitative Summary}

\subsection{555897: How can the prep material be improved?}

By prep material we mean the papers to read and annotate in FeedbackFruits and 
the short reflection essays that you submitted in Canvas before each seminar.

\subsubsection{Individual Responses (40 total):}

\begin{enumerate}
  \item I think some of the papers might have been obsolete/very old, I also think it would have been nice to have reading material aimed more at methodology, perhaps specifically written for the course.
  \item I did the project instead of seminars
  \item It was good
  \item It was good
  \item Fewer papers.
  \item I missed the first seminar so I focused on the exam.
  \item Feedbackfruits is very annoying to use so would be good if it could be done in some other way.
  \item I think the preparation material is great.
  \item The biggest problem with prep material is that it wasn't released early enough. We didn't have enough time to go through it all and plan ahead before the seminars.
  \item I didn't do the reflection part before the seminar, but the papers could be improved by having some kind of correct answer, where the teacher highlights the parts where the information can be found.
  \item Didn’t read it, did the final assignment 
  \item The biggest area of improvement isn't within the material itself but the uploading of the material. Rather than uploading a week before the seminar (and in some cases days before the seminar), it would have been nice if they could be uploaded earlier. By doing so, we as students can more easily find slots within our schedule where we can read and prepare. 
  \item In my opinion, there were way too many prep material for each seminar and they were released way too close to the actual seminar which resulted in a very rushed preparation. I think a better solution is to have less papers to read. I think when there are as many papers as there were in this course, a lot of students just try to rush through them which makes it difficult to learn. Also along these lines, I think a lot of papers just gives an unnecessary workload, I think the concept this course is trying to introduce regarding finding the research question and methodology etc is well introduced by one or two papers and I think more papers just wastes time because the student is stuck reading lengthy papers in order to satisfy the criteria for passing.   In summary, less papers to read for each seminar, and release the papers earlier prior to the seminars.
  \item I found the FeedFruit interface confusing more than helping. Just reading the article and then answering what was research question, research methods used separately would fit me more.
  \item They were good 
  \item (Video) lectures explaining concepts and giving short examples would be preferable to workload-intensive papers
  \item They are a bit too long, maybe you could find a better version that summarizes the paper down a bit but still has all the needed information
  \item The prep material is good
  \item Kind of make sure that actually everybody read the papers beforehand. Quite often I was in breakout rooms and the minority of people actually read the papers...
  \item I believe that the prep material could be improved by having less technical papers and more focused on the methodologies that we were discussing. Also I believe that some of the annotations that we had to do on the papers were poorly explained on how we had to annotate exactly, more examples would have been nice.
  \item I think the papers were too big. Most of them were difficult to understand. 
  \item Especially for the papers it would be nice to have at least one example or maybe some kind of possible solution (i.e. some feedback on how correct the annotations are). Even now I am not sure if I annotated correctly and sometimes I wasn't even sure what use the annotations were because I did not know if I got it wrong (especially for the green annotation - I think that was methodology - I have no clue if I'm on the correct path and if that has some connection to Tamos, because lets be honest the papers do not say "we did a lab experiment with bliding and randomization", which makes that whole course super abstract).
  \item Papers on computer and cybersecurity cover a wide range of topics. If a person does not have the knowledge background of related technologies, reading them will make people lose confidence and interest. For example, I find it more enjoyable to read papers on password and cryptography, but it is more difficult for me to read papers on onion routing because I have little knowledge of this area. In addition, it is easy to get lost in the technical details when reading these technical papers, which deviates from the goal of this course.  I think more specific reading guidance and reflection questions can improve this situation. Multiple-choice quiz that reveals the answer after completing them  could be used to help everyone clarify the reading goals and directions. Complete demonstration example on how to analyze the methodology for a paper could be helpful for students to learn. It may also be helpful to point out which parts of the paper can be ignored.   Some directly related reading materials can allow students to have some understanding of relevant knowledge before reading a paper in an unfamiliar field, and increase their confidence in reading papers.
  \item I feel that the material that needs to be read can be identified by marking the key points as well as the core parts of the paper, which will increase the efficiency of reading the paper, and the time saved can even be used to go and read a few more papers.
  \item I am happy with the prep material. I felt I could discuss during the seminars frequently and that I understood the subjects. I felt like I could use concepts from Tamos even though I studied it during P1. It was good repetition. 
  \item I think getting these assignments earlier would help a lot since sometimes we got most of the papers  1 - 2 days before the seminar which made it more of a stress run where at least myself wouldn't have time to contemplate or consider what I had read. 
  \item I only did the final assignment.
  \item The consistency of the FeedbackFruits activities can definitely be improved (sometimes the reflection in the FBF activity was the default one, sometimes custom, sometimes you didn't have to write a summary for a section, etc), as well as the clarity in what we're expected to write and reflect about more precisely than just "research question". 
  \item In my opinion, there were too many papers to read which made it difficult to remember the content of all of them for the seminar. It was also difficult to know which assignments were necessary for the seminar and which weren't.    
  \item Sometimes there were too many papers but I think this was mainly because they were sometimes posted a bit late which is understandable because the course was still undergoing changes. However, most papers were interesting and helped with the reflections by brining up points that I would not have thought of. I think the reflections were the best as prep material as it was mainly the reflection that was discussed during the seminars. 
  \item Feedback fruits feels very clunky at displaying information. You have to navigate around menus looking for information which might not exist, it generally felt pretty confusing using it. 
  \item I thought the structure of the course was really unclear. I was prepared to do all the participation assignments and seminars instead of the project, but I felt like I did not what assignments to do before each seminar. It did not help that the information and assignments were released just within a few days as there are other courses taking a lot of time.
  \item I think some papers are quite complex in terms of the associated languages they have. Maybe you could highlight what one should focus on in each of those papers in order to not "go on a tangent" of their own or simply chose a smaller set of these papers
  \item The prep material could be improved by having a broader set (different) of questions regarding the papers other than what was the research question, method etc.
  \item At times the prep material was released a bit too late. Because we couldn't see the next seminar and its prep material until you published it, which would be within one week, often a half, would mean that you had to all of a sudden put in a lot more time than you had anticipated during those days/week. I plan my workload throughout the week based on information the courses provides, so when I don't have information when I plan and then suddenly I have papers I need to read and annotate, or write a short reflection essay, then my plan doesn't work anymore. Then I have to readjust my plan (for the other courses) which at times was hard depending on upcoming deadlines.   In short, I would appreciated it information was released sooner so that I could plan my days ahead of time easier. I only attended the last mandatory seminar, partly because of this, but still tried to do some prep work thinking it would still benefit me.
  \item It was a lot of technicalities in the papers and they were so long too, they need to be shorter if they are that difficult to read.
  \item Shorter articles. Some examples of what you want us to annotate perhaps.
  \item i like feedback fruits in general, good tool but integration to canvas is somewhat lacking. i do not have any more specific improvement suggestions
  \item I think perhaps the prep material is good, but I feel like when we've read it we just left them in the dust kind of thing. I would like to discuss them further in the seminars to make sure that e.g. we indeed found the correct research question which was required to annotate in feedbackfruits to get a passing grade.   Moreover, the prep material has to have a more clear structure. It was very confusing knowing what had to be done and what was optional (this also changed during the course).
  \item For me, more videos, less papers.
\end{enumerate}

\subsubsection{AI-Generated Summary:}

\begin{itemize}
\item \textbf{Main themes}
  \begin{itemize}
    \item \textit{Workload and timing:} Many felt there were too many and/or too long papers, often released too close to the seminars, making preparation rushed and hard to plan.
    \item \textit{Fit to course goals:} Several papers were perceived as overly technical or hard to understand; students wanted materials more focused on research methodology and tailored to the course (some flagged papers as old/obsolete).
    \item \textit{Tooling and process:} FeedbackFruits (FBF) was frequently described as clunky/confusing, with inconsistent activity setups and unclear expectations for annotations.
    \item \textit{Guidance and feedback:} Students sought clearer instructions, exemplars, and “answer keys” for annotations; more targeted prompts, highlights of key sections, and help distinguishing what to focus on.
    \item \textit{Format variety and seminar connection:} Interest in shorter readings, summaries, and short video lectures; desire to use seminar time to validate findings (e.g., research questions) and ensure alignment.
  \end{itemize}

\item \textbf{Common concerns or issues}
  \begin{itemize}
    \item Late release of prep materials and insufficient lead time to plan (\textit{mentioned repeatedly}).
    \item Excessive number/length of papers; difficulty retaining many papers’ content for discussion.
    \item Papers too technical or outside students’ background; hard to connect details to methodology.
    \item FeedbackFruits usability issues; confusing navigation and Canvas integration; inconsistent templates and unclear annotation instructions.
    \item Unclear course structure and requirements (what is required vs. optional; changing instructions).
    \item Limited feedback loop on whether annotations (e.g., research question, methodology) were “correct.”
  \end{itemize}

\item \textbf{Suggestions for improvement}
  \begin{itemize}
    \item \textit{Scheduling and scope:} Release materials earlier; reduce the number of papers per seminar (e.g., 1–2 core papers); prefer shorter/less technical papers or provide condensed versions.
    \item \textit{Methodology alignment:} Choose or create course-specific methodology guides; ensure papers support methodological learning over technical deep dives.
    \item \textit{Guidance and support:} Provide exemplars/sample annotations, rubrics, and brief “solution keys”; highlight key sections and what can be skimmed; offer background primers for unfamiliar topics; include focused reading goals.
    \item \textit{Activities and assessment:} Add short auto-graded quizzes (with revealed answers) to clarify objectives; diversify reflection prompts beyond “research question/method”; pre-seminar checks to encourage completion.
    \item \textit{Platform and instructions:} Streamline or reconsider FBF usage; improve Canvas integration; standardize FBF activity templates; present a clear per-seminar checklist labeling required vs. optional tasks.
    \item \textit{Seminar use:} Dedicate time to collectively confirm research questions/methods and discuss key takeaways from the readings.
    \item \textit{Format variety:} Incorporate short video micro-lectures or concept explainers to reduce reading load and reinforce methodology.
  \end{itemize}

\item \textbf{Overall sentiment}
  \begin{itemize}
    \item Mixed. A notable minority praised the materials as good/useful, but the prevailing concerns centered on logistics (late release, unclear structure), workload (too many/long/technical papers), and tooling (FBF usability).
    \item Students generally value reading-based preparation when it is lighter, released earlier, better guided toward methodology, and paired with clearer instructions and feedback.
  \end{itemize}
\end{itemize}

\subsection{555899: How can the seminars be improved?}

For example the topics discussed, the grouping into breakout rooms, shorter or 
longer times for discussions, longer seminars to cover the topics more deeply.

\subsubsection{Individual Responses (38 total):}

\begin{enumerate}
  \item I did the final project and thus not the seminars, but did look through the papers.
  \item I did the project instead of seminars. 
  \item It was good
  \item The seminars where really well conducted
  \item I would like to meet in person.
  \item I don't have anything to add.
  \item It would be great if there is a list of seminars to attend in the home page so we don't miss the initial ones.
  \item They're fine
  \item I didn't attend the seminars, so I can't judge
  \item I liked the final assignment seminar, can’t speak for the others
  \item I think the seminars were good, I can't really think of anything that required improvement. 
  \item I think the seminars are good as they are, more time to prepare before each seminar would be helpful to ensure a deeper understanding of the subject before participation.
  \item I would held the seminars in person. I just feel there are always some glitches, someone mic does not work and it's just feels ineffective.
  \item By dividing the seminar topics or research paper into smaller groups so that the whole class can then discuss about what the groups found 
  \item Longer times for discussion
  \item In my opinion the seminars were great, it is mostly great that students get the chance to discuss within them self's and find new information that they might have missed.
  \item Larger breakout groups would encourage more active participation, and extending discussion time for complex topics like side-channel attacks could enhance understanding.
  \item Breakout rooms worked from time to time very well, but then again other times not. I was partially stuck in a breakout room where nobody responded to my questions, or the discussion was completly in Swedish (I don't speak any) or the people were not very well prepared...
  \item I believed the seminars were good, I think longer time for discussions is necessary, I believe the seminars would benefit the students especially if the discussions in breakout rooms were longer and better guided, with questions that are more open in order to discuss them for a long time.
  \item Maybe in the seminars the teacher could explain us some topics, I prefer this type of learning. 
  \item I think this was done really well :-) .
  \item I remember that during the last discussion, the groups were not balanced enough. I felt that the discussion in a group with too many people was not as good as that in a group with fewer people.   In addition, I like the discussion to be well-founded rather than general. This may require us to narrow the scope of the discussion, or the issues discussed are specific and precise, telling us which materials or angles to think about. And it may be better for everyone to think about takeaways after the seminar.  By the way, shared document editing which was used in the seminars is a good way, where everyone can write down their own ideas and examine themselves through words :)
  \item I don't think there's anything that could be done to improve the seminar, it's great as it is!
  \item Longer seminars I think. It felt like I would have easily been able to discuss the subjects for another hour or so. I liked them a lot. 
  \item In general, the seminars were good, I think they helped me understand many of the topics better by discussing and hearing others' opinions on these topics. One way to improve these would be to allocate more time to the seminars or maybe divide the module up and allocate more seminars since we couldn't go through everything on the seminars. 
  \item No opinion; I only did the final assignment.
  \item More consistent placement, having one seminar per week instead of two seminars every other week would be better in terms of making the workload more even.   This also ties in with the prep work, but releasing all prep work for the seminars in good time (like 1 week in advance) and being clear with what needs to be done for each seminar would also be appreciated. 
  \item I think it would be better if we only focused on a few of the papers that we read, that way it's easier for everyone to remember what was said in those reports. I also think it would be better if we got to choose our own groups because it's more likely that we'll talk more if it's people that we're comfortable with.
  \item I think the seminars were good. The topics reflected the papers and the reflections. The discussion times were adequate except for the last seminar where the time was a bit short to discuss all the chosen questions which we designed  experiments to. 
  \item No comment 
  \item I think the areas talked about in the seminars were good, however, the material or slides for the seminar could be released beforehand. It felt like the material on the seminars differed a lot from the prep material and it would have been nice to know what to focus on more to make the seminars effective.
  \item * groups were fine  * perhaps a better solution than using Google Docs could work in the future, but then again, that would take more time to set  * Use through-provoking language in the topics might increase the enthusiasm
  \item I think that the seminars had a good structure good. The information and prepp assignments where released close to the seminars which made us student unable to create a study plan for the course. 
  \item I only attended the last mandatory seminar for the project. How that seminar could have been improved is mostly information about seminar expectations. Maybe I would have had a better idea on what to expect if I had attended the others. The seminar was good, but I went into it thinking we had to present and discuss all research questions we found or talked about in the report. But we (my group) quickly realized that we only needed to present our most interested findings and discuss that one so we quickly pivoted. Maybe this information was available somewhere but we had some general confusion about the instructions throughout the project. 
  \item It was okay as it was
  \item i like the time for each seminar, it was unfortunate thet one seminar got canseld. I like the discussions that was good. also the presentation by the kth libary was nice. 
  \item Fewer but longer seminars to discuss the papers read (specifically the methodology) to gain a deep understanding of the methods that are commonly used in cyber security. Briefly skimming through an analysis of a cryptographic proof  and then discussing it for 5 minutes makes it difficult to comprehend. 
  \item In person maybe? But I liked zoom also.
\end{enumerate}

\subsubsection{AI-Generated Summary:}

\textbf{Main themes}
\begin{itemize}
\item Broad satisfaction with seminar quality and usefulness; several respondents did not attend (project track or only final seminar).
\item Strong interest in more time and depth for discussions, especially on complex topics and methodologies.
\item Breakout-room experience is pivotal: group size/balance, participant preparedness, and language consistency affect quality.
\item Modality and logistics matter: some prefer in-person to avoid tech glitches and improve engagement.
\item Scheduling and preparation: desire for earlier release of prep work, slides, and a clear seminar schedule/cadence.
\item Structure and facilitation: requests for clearer expectations, focused scopes, and guided yet open-ended questions; occasional teacher mini-explanations desired.
\item Tools: shared collaborative documents appreciated; some skepticism about Google Docs and interest in alternatives.
\end{itemize}

\textbf{Common concerns or issues}
\begin{itemize}
\item Insufficient discussion time in some seminars (notably the last one).
\item Uneven breakout rooms: oversized groups, unbalanced participation, silence, language issues (non-English), and variable preparedness.
\item Prep materials/slides released too close to seminar; mismatch at times between prep content and seminar focus.
\item Irregular scheduling/cadence and a canceled session increased planning difficulty.
\item Occasional Zoom/audio glitches and uncertainty about seminar expectations (especially for the project seminar).
\end{itemize}

\textbf{Suggestions for improvement}
\begin{itemize}
\item Increase time for discussion: longer seminars or fewer-but-longer sessions; allocate extra time for intricate topics (e.g., side-channel attacks, proofs/methods).
\item Optimize breakout rooms: balance group sizes, ensure language consistency, encourage preparedness; consider topic-based subgroups that report back; allow group choice in some cases.
\item Strengthen planning and clarity: publish a full seminar list on the homepage, adopt a steadier cadence (e.g., weekly), release all prep work/slides at least a week in advance, and state expectations/instructions clearly.
\item Focus and depth: narrow scope, prioritize fewer papers per session, use specific prompts and thought-provoking language; end with explicit takeaways.
\item Facilitation: include brief instructor explanations to anchor complex material; provide more open, well-guided discussion questions.
\item Modality and tools: offer in-person or hybrid options where possible; keep shared-doc collaboration but evaluate more robust alternatives to Google Docs.
\end{itemize}

\textbf{Overall sentiment}
\begin{itemize}
\item \textit{Mostly positive to very positive.} Many found seminars well-conducted, engaging, and helpful for understanding through discussion; feedback centers on logistics, time allocation, and improving breakout dynamics rather than on core design.
\end{itemize}

\subsection{555900: What was the best part of the seminars and prep?}

For example a favourite paper, an eye-opening topic, or just to discuss how 
this relates to TAMOS.

\subsubsection{Individual Responses (38 total):}

\begin{enumerate}
  \item I did the project instead of seminars
  \item Discussions
  \item .
  \item Finally understand why TAMOS is important.
  \item Didn't attend seminars
  \item It is that the papers were helpful during the TAMOS seminars to discuss from the security perspective
  \item Haven't done enough prep to answer this
  \item For me, being able to see the the applications of TAMOS from a Cyber security point of view was very interesting, to see how a cybersecurity scientist works and does research.
  \item .
  \item The best part was the seminar itself where you got the upportunity to talk with other students and understand their perspective of the respective subjects. 
  \item The best part was to utilise the things we learned in the TAMOS course to analyze papers. 
  \item Dos and Don'ts of Machine Learning in Computer Security 
  \item it was interesting to discuss and see how tamos lecture be used in security 
  \item Matching TaMoS to Cybersec
  \item I really liked every part of the seminar, it gave a more thorough insight on the papers we read and it was fun to see that others also had the same opinion.
  \item Discussing how TAMOS methodologies apply to security problems was a highlight. The paper on ElGamal encryption and side-channel vulnerabilities was especially thought-provoking.
  \item The ongoing discussions with you and the whole group. It was interesting to see the different approaches and interpretations for different topics.
  \item For me the best part was discussing all of the ideas in a practical manner with other students, since that helped me understand the topics more deeply.
  \item The discussion in the seminars were the best part. 
  \item I found especially eye-opening the first video (where you explained a few basic tamos concept with security examples) and I really liked the discussions together (everyone) in the seminars.
  \item There are two papers that impressed me deeply. One is a paper on password, which made me realize that investigating user behavior is a qualitative method in cyber security, which supplements the content that cannot be revealed by technical quantitative methods. The other is the Elgamal paper. In the cryptography course I took before, I was exposed to cryptographic algorithms from the perspective of mathematical concepts and proofs, while that paper explained that there is an inconsistency between the principle and implementation of the algorithm. Even if the algorithm is secure, the implementation may still be insecure.
  \item I can read a lot of papers in preparation for the seminar, which will help me learn a lot of things
  \item I liked it when we had to write methodologies for different problems. That was very fun and the discussion we had in one of the groups was also great. That was probably one of my favorite parts of the course. Perhaps there can be a division so that groups are small sizes since the password group contained a lot of people. I wasn't in that one but it must have been messy. 
  \item For me the most interesting topic was on the last one where we were supposed to come up with methods for making sure something is secure. This allowed me to apply the theory better than what we did on the seminars in TAMOS since TAMOS is a more general course. 
  \item No opinion; I only did the final assignment.
  \item I definitely enjoyed the part about formal methods and model checking the most. It's something I've been semi-interested in even before this course, and it was interesting to see how it could be discussed and used in relation to security (apart from just "haha program work correctly") as a method.
  \item I think the content of the papers were really interesting and it gave new perspectives that the TAMOS course didn't include.
  \item My favourite paper was the SoK Science, security and the Elusive goal of security as a scientific pursuit as it gave a clear understanding of the difficulty of proving that something is secure and that you have to find other ways to show that something is secure "enough". I also like the passwords and models modules in general because there were some interesting papers were cool studies done. 
  \item No comment 
  \item The best part was finding correlations between security and TAMOS, considering how abstract the material can get in philosophy. 
  \item I liked the diversity of topics although some were difficult to follow up on a first viewing
  \item It was interesting to discuss with other people to get different views of topics.
  \item As I said, I only attended the project seminar and did some of the prep work so I don't think I have anything really to add to this. 
  \item I liked discussing with my classmates
  \item I enjoyed the password-papers.
  \item i think this course for it provides acuall understanding and usfulness of TAMOS to practical real topics for me and other CS students. I think TAMOS have forgotten about us and do not understand our examples
  \item Getting some picked papers to read which conduct a sound and complete methodology in cyber security was helpful for future projects. Moreover, the fact that some cryptographic protocols or algorithms can be broken given a specific context was interesting!
  \item Reflections and videos.
\end{enumerate}

\subsubsection{AI-Generated Summary:}

\textbf{Main themes}
\begin{itemize}
\item \textit{Discussion-centric learning:} Peer discussions and hearing diverse perspectives were repeatedly cited as the best part.
\item \textit{Bridging TAMOS to cybersecurity:} Applying TAMOS methodologies to concrete security problems helped clarify why TAMOS matters and made abstract ideas feel practical.
\item \textit{Paper-driven insights:} Notable favorites included:
  \begin{itemize}
  \item Passwords and user behavior (value of qualitative methods alongside quantitative).
  \item ElGamal and side-channel/implementation gaps (principle vs. practice).
  \item SoK “Science, Security and the Elusive Goal of Security as a Scientific Pursuit.”
  \item Formal methods and model checking.
  \item “Dos and Don’ts of Machine Learning in Computer Security.”
  \end{itemize}
\item \textit{Methodology practice:} Analyzing papers with TAMOS tools and designing methods for security questions were engaging and instructive.
\item \textit{Support materials:} Introductory videos and reflections that tied TAMOS concepts to security were appreciated.
\item \textit{Breadth:} Diversity of topics was valued, though sometimes challenging on first pass.
\end{itemize}

\textbf{Common concerns or issues}
\begin{itemize}
\item Some students did not attend or had limited preparation, limiting their ability to comment.
\item Large group sizes (e.g., the password group) made discussions messy and reduced individual participation.
\item Certain materials were hard to follow on first viewing; TAMOS can feel abstract or misaligned with CS students without concrete examples.
\end{itemize}

\textbf{Suggestions for improvement}
\begin{itemize}
\item Use smaller seminar/breakout groups with capped sizes to improve engagement.
\item Provide more scaffolding: brief primers, step-by-step methodology templates, and security-focused examples linking TAMOS concepts to practice.
\item Curate papers with clear, sound methodologies and real-world implementation pitfalls; balance difficulty with accessibility.
\item Include more hands-on activities where students craft and critique methods; add short prep guides or guiding questions to encourage consistent preparation.
\end{itemize}

\textbf{Overall sentiment}
\begin{itemize}
\item Predominantly positive among participants. Discussions, practical application to security, and selected papers were highlights that deepened understanding and demonstrated TAMOS’s relevance. Issues raised were mainly logistical (group size, pacing/clarity) rather than content quality.
\end{itemize}

\subsection{556032: How can the project assignment be improved?}

\subsubsection{Individual Responses (31 total):}

\begin{enumerate}
  \item I think the instructions said that we should have come up with "all" research questions or something like that, I believe that is (almost) impossible. Instead it should be one, or two, or 5 or some specified number.
  \item The instructions were somewhat vague, particularly regarding the scope and the expected time commitment. We couldn't find any clear guidelines, such as page or word limits which led to much of the project not being included in the final seminar. Additionally, the project was published on December 20th, and an earlier release would have been preferable to avoid working on it during Christmas.
  \item It would be good to inow at tthestart of the course if the seminars are enough to pass the course or not
  \item It was good
  \item N/A
  \item I was a bit confused about the instructions. Would be nice to have an example.
  \item I think there should be more topics to choose from
  \item Some information about the project was only clear through the slides shown during the presentations seminar. It was also difficult to understand what we were supposed to present in the seminar and we ended up making a lot more slides than necessary.
  \item The overall project was very well structured, making the students reflect on their knowledge and design themselves a methodology to answer a question. I wouldn't change anything from the final assignment.
  \item More clearly instructions, specifically on the presentation 
  \item IDK
  \item The guidelines could be more precise as to what to expect from the project. I found some major differences in the approach that my group took and the group I reviewed. Both out documents were similar but some key differences arised when designing a methodology.
  \item The project assignment was clear and reflected the course well.
  \item -
  \item By putting out there the instructions to project so it is easier to plan ahead if one wants to do it 
  \item Introduce the requirements and expectations better before the project, not during the presentation seminar at the end
  \item I didnt need to do the project.
  \item No project assignment
  \item I didn’t have to do the project.
  \item Thank you!
  \item I don't have any idea for this.
  \item At the project allocation stage, the allocation of each group in terms of how many people should be made clear, as well as the structure of the project report.
  \item I passed without the project
  \item No idea since I didn't need to do it.
  \item The project assignment itself was mostly fine in my opinion, but the seminar structure was a bit unclear at first and I didn't know what to expect of it.
  \item While I didn't do the project assignment, it would've been nice if information regarding the project assignment was released earlier in the course. 
  \item The information about the project felt too sparse and feedback fruits was confusing to use. For example you had some points for "adjustment" in Canvas, but the feedback fruits activity didn't feature anywhere to turn in any adjusted text. This could either mean that the adjustment wasn't necessary (outdated instruction) or the task wasn't posted yet, or we had missed something vital, or it meant adjusting it for the seminar presentation.   Another example is the early canvas deadline for peer review of the text (which no one had even had the option to turn in yet). It was just generally confusing and it felt like you had to scavenge for information between canvas and hidden feedback fruits pages.  Another example is going to activities within feedback fruits. It simply loads infinitely, which left you wondering if there were no activities, no released activities or a bug in feedback fruits. 
  \item Did not do
  \item Even though the instructions contained a lot of information we had trouble figuring out the expectations we had to meet. The biggest thing was the project seminar which I mentioned in an earlier question, but also when and how we were going to submit the final version of the report. This was because the "final report" assignment was unpublished and hidden from us, and we only saw it during the report when Daniel mentioned it. We were debating between us where we could submit the final, and if we even needed to do a final version or if the seminar counted as that. This and the fact that deadlines was between Christmas and New Years made it really difficult to find time to work, work with group members, and communicate with Daniel. If I had known that the deadlines was during such busy days I would have made the extra effort to do the seminars. I also point this to content not being published until late, as we had no info on the project until december, and at that point it was too late for me to do the seminars as I already had missed one or more.
  \item some more clarity in the instructions and specificly on how the slides shuld be. Also some proof read of the instruction to remove old info relating to another version of the final seminar. It relly sucked to miss one seminar and have to do the project... maybe allow one missed seminar in the future. but the final seminar was good, discussions and info
  \item First off, define clearly in the beginning of the course if the seminars are optional or not. In the introductory lecture they were stated as optional which was later decided after seminar 1 that if you attend all seminars you get a passing grade. I think this should have been communicated both verbally and through an announcement on canvas before seminar 1.   Secondly, perhaps we over-interpreted the project assignment but when it says "All existing literature", we took that literally and spend all of christmas break reading scientific literature on the topic of Signal which was a large amount of work to do. Only to then being with the actual assignment. A rephrasing of the assignment would be helpful.
\end{enumerate}

\subsubsection{AI-Generated Summary:}

\textbf{Main themes}
\begin{itemize}
\item \textit{Clarity and consistency of instructions}: Frequent confusion about scope, expected workload, deliverables, and presentation requirements.
\item \textit{Timing and communication}: Late release (Dec 20), holiday deadlines, and shifting/unclear rules (e.g., whether seminars alone suffice to pass).
\item \textit{Seminar expectations}: Unclear what to present, how many slides, and how seminars relate to grading.
\item \textit{Tools and submission flow}: Fragmented information between Canvas and FeedbackFruits; unpublished/hidden assignments; peer-review timing mismatches; platform glitches.
\item \textit{Scope/workload definition}: Ambiguous phrasing (e.g., “all research questions,” “all existing literature”) led to over-interpretation and unnecessary workload.
\item \textit{Logistics}: Group size allocations and report structure not clearly specified; limited topic choices.
\item \textit{Positives}: Several found the project well structured and the final seminar valuable.
\end{itemize}

\textbf{Common concerns or issues}
\begin{itemize}
\item Vague guidelines (no page/word limits, unclear time commitment, differing interpretations of methodology).
\item Unclear seminar format and presentation expectations; overproduction of slides.
\item Late publication and holiday deadlines hindered coordination and effort.
\item Inconsistent communication about pass pathways (seminars vs. project) and unpublished submission points.
\item Confusing, sometimes outdated instructions; conflicting information across platforms; FeedbackFruits loading/visibility issues.
\item Overly broad wording (“all literature,” “all research questions”) felt unrealistic.
\end{itemize}

\textbf{Suggestions for improvement}
\begin{itemize}
\item Specify concrete scope: exact number of research questions, page/word limits, expected time-on-task, and clear methodology expectations.
\item Clarify seminar role and grading policy at course start; state explicitly whether attending all seminars yields a pass.
\item Release project information and all related assignments earlier; avoid holiday deadlines; publish a clear timeline.
\item Provide exemplars/templates (report and slides) and explicit presentation guidelines.
\item Consolidate instructions on one platform; ensure all tasks are published and visible; fix FeedbackFruits access/performance; align peer-review and submission dates.
\item Proofread to remove outdated references (e.g., “adjustments”) and rephrase unrealistic terms (“all existing literature” → “representative/key literature”).
\item State group sizes and report structure up front; offer more topic options.
\item Consider allowing one missed seminar without defaulting to the full project.
\end{itemize}

\textbf{Overall sentiment}
\begin{itemize}
\item Mixed: Many appreciated the project’s intent and structure, but experienced significant confusion due to unclear, late, and sometimes conflicting instructions and logistics. A notable share did not complete the project, and their feedback focused on timing and communication rather than content.
\end{itemize}

\subsection{564845: Do you have anything else you would like to add to the evaluation of the course?}

\subsubsection{Individual Responses (15 total):}

\begin{enumerate}
  \item It wasn't communicated at the start of the course whether the seminars alone were enough to pass, so completing the project seemed like the only option. It would be better if this were clearly stated from the beginning.
  \item No
  \item No, thank you!
  \item Is this quiz the reason for why my grade is not in Ladok yet?
  \item If you missed one seminar it felt like there wasn't any point to attend the others and just focus on the project instead.
  \item I hated the course! Just kidding, but I don't think it was very useful. I think 4,5hp + 3hp is too many credits taken from other things. I think TAMOS was a very good course (good structure) and I learned a lot, so I would not want that taken away. That leaves this course. I think this course could be 0,5 hp and it wouldn't make any difference. Maybe re-integrate it into cybsam.
  \item No nothing more
  \item The course felt rushed with the seminar preparation assignments being released very close to the actual deadlines, making it almost impossible to do everything on time and plan ahead.  Additionally, I would have preferred to do the seminars instead of the project, but only if I had certainty I would pass the course with that. Hopefully something that's here to stay in the next editions of the course.  Finally, I took TAMOS before this course (P1 offering was cancelled), so I was a bit disconnected since I was no longer engaged in the course subject.
  \item No
  \item No
  \item NO
  \item no.
  \item I found the organization of the course little confusing. Not sure exactly why, but I was not sure what I am suppose to do until the first seminar.
  \item No
  \item It was very confusing to know whether active participation would be counted or not to pass the course, so ultimately I had to prioritize other courses and had to do the project. If I had known from the start, I would have made a deeper effort to do the prep work and attend all the seminars (except the one that clashed with the Forensics guest lecture -- it would be good for seminars to be listed in the calendar from the get-go to prevent this).
\end{enumerate}

\subsubsection{AI-Generated Summary:}

\textbf{Main themes}
\begin{itemize}
\item Clarity of pass requirements and assessment pathways (seminars-only vs. project; whether active participation counts).
\item Course organization and scheduling (late release of seminar prep, difficulty planning ahead, calendar visibility, clashes with other activities).
\item Attendance incentives and motivation (missing one seminar undermined motivation to attend the rest).
\item Perceived workload/credit value (course seen as low value for its credits; comparison to TAMOS).
\item Administrative transparency (uncertainty about grade posting in Ladok).
\end{itemize}

\textbf{Common concerns or issues}
\begin{itemize}
\item Unclear from the start whether seminars alone were sufficient to pass; confusion about requirements until the first seminar.
\item Seminar preparation released too close to deadlines, creating a rushed feel and planning challenges.
\item Seminar times not listed early, causing conflicts (e.g., Forensics guest lecture).
\item “All-or-nothing” seminar policy discouraged continued attendance after missing one.
\item Question about whether a quiz delayed grades being posted to Ladok.
\item Sequencing/cancellation effects (taking TAMOS earlier led to disconnect from this course).
\end{itemize}

\textbf{Suggestions for improvement}
\begin{itemize}
\item State pass criteria and assessment options clearly from day one (e.g., whether a seminars-only path is possible and the exact requirements).
\item Publish all seminar dates and add them to the calendar at the start; avoid clashes with core/guest lectures.
\item Release seminar prep materials earlier with realistic deadlines to support planning.
\item Provide flexibility after a missed seminar (make-up options or thresholds) to maintain incentives to attend.
\item Communicate administrative timelines (e.g., Ladok posting, impact of quizzes) proactively.
\item Reassess credit load and alignment; consider reducing credits or integrating into a related course (while keeping TAMOS intact).
\item Coordinate timing with TAMOS offerings to minimize disconnect when schedules shift or are cancelled.
\end{itemize}

\textbf{Overall sentiment}
\begin{itemize}
\item Mixed, leaning negative among substantive respondents: primary dissatisfaction centers on unclear requirements and rushed scheduling.
\item Many responses were neutral/no additional comments (8/15), but detailed feedback highlights confusion, planning difficulties, and questions about the course’s value relative to its credits; several would prefer the seminar path if assured it leads to passing.
\end{itemize}


\chapter{vetcyb25p4}

\section{Quantitative Summary}

\subsection{564848: How much time did you spend on the teaching of the course?}

That is until the teaching ended before the exam period. (The expected workload 
is 80 hours for the whole course, so roughly 8 h/week for 10 weeks. The 
teaching period is 7 weeks.)
Don't include the time for watching the TAMOS lectures, I count that time to 
the TAMOS course.

Total responses: 23\\

\subsubsection{Response distribution:}
\begin{itemize}
  \item 6-8 h/week: 9 (39.1\%)
  \item 2-4 h/week: 5 (21.7\%)
  \item 4-6 h/week: 4 (17.4\%)
  \item 10-12 h/week: 2 (8.7\%)
  \item 8-10 h/week: 2 (8.7\%)
  \item 0-2 h/week: 1 (4.3\%)
\end{itemize}

\subsection{564849: Below you'll find a few statements about the prep material. Check all that you consider true.}

Total responses: 23\\
Total selections: 54\\

\subsubsection{Option distribution:}
\begin{itemize}
  \item It was helpful to see others' annotations in the papers that we had to read.: 15 (65.2\%)
  \item It was a suitable number of papers to read for each seminar.: 11 (47.8\%)
  \item It's a good idea to annotate the TAMOS lectures from a security perspective.: 11 (47.8\%)
  \item It was too many papers to read for each seminar.: 9 (39.1\%)
  \item Some topics would have benefited from more papers as example methodologies.: 3 (13.0\%)
  \item Annotating the TAMOS lectures would be more useful with more guidance.: 3 (13.0\%)
  \item It was too few papers to read for each seminar.: 1 (4.3\%)
  \item It would be helpful to not see the others' annotations in the papers until done with one's own annotations.: 1 (4.3\%)
\end{itemize}

\subsection{564850: Below you'll find a few statements about the prep material.Check all that you consider true.}

Total responses: 23\\
Total selections: 41\\

\subsubsection{Option distribution:}
\begin{itemize}
  \item Reading the papers was good prep for the seminars.: 17 (73.9\%)
  \item Writing the reflections was good prep for the seminars.: 13 (56.5\%)
  \item The reflection instructions can be improved to better prepare us for the seminars.: 9 (39.1\%)
  \item Some topics should have had more papers to better prepare for the seminar.: 2 (8.7\%)
\end{itemize}

\subsection{564852: Below you'll find a few statements about the seminars. Check all that you consider true.}

Total responses: 23\\
Total selections: 50\\

\subsubsection{Option distribution:}
\begin{itemize}
  \item It was helpful to discuss the topics with classmates during the seminars.: 18 (78.3\%)
  \item The seminars helped me deepen my understanding for the science of security.: 15 (65.2\%)
  \item The seminars helped me better understand how the methods learned in TAMOS can be used in security.: 14 (60.9\%)
  \item Some that I discussed with during a seminar weren't prepared so they couldn't contribute anything to the discussion.: 3 (13.0\%)
\end{itemize}

\subsection{564855: Below you'll find a few statements about the project.Check 
all that you consider true.}

Total responses: 23\\

\subsubsection{Response distribution:}
\begin{itemize}
  \item I didn't have to do the project since I did all prep and participated : 20 (87.0\%)
  \item I missed a seminar so I had to do the project.: 2 (8.7\%)
  \item I never planned on participating during the course\, I was aimin for t: 1 (4.3\%)
\end{itemize}

\subsection{571034: I want to use the data in the course for pedagogical 
research and development.}

To do so, I need your consent. You can always opt out again by updating your 
answer here or emailing dbosk@kth.se.
In any research work, I will remove all references to your identity from the 
data (it will be \enquote{anonymized}; don't use steganography to hide your 
identity in the data, I can't find and remove that ;-)).

Total responses: 23\\
Total selections: 46\\

\subsubsection{Option distribution:}
\begin{itemize}
  \item Yes: 17 (73.9\%)
  \item I agree that my data can be used in research.: 17 (73.9\%)
  \item No: 6 (26.1\%)
  \item I don't want my data to be included in research.: 6 (26.1\%)
\end{itemize}

\section{Qualitative Summary}

\subsection{564851: How can the prep material be improved?}

By prep material we mean the papers to read and annotate in FeedbackFruits and 
the short reflection essays that you submitted in Canvas before each seminar.

\subsubsection{Individual Responses (22 total):}

\begin{enumerate}
  \item I do not know actually, since I found the prep material already very well presented! Sorry for not being of much help!
  \item The prep material was pretty good. 
  \item More explanations about key concepts
  \item One way to improve could be to give more options in terms of papers to read or allow us students to find relevant papers ourselves before some of the seminars. By doing so, the seminars could benefit from added perspectives which could add to the discussions.
  \item I think most people did not read through all the material as there was 1. a lot of papers to read and 2. you did not have to read all of the paper to complete the tasks most of the time.
  \item I believe FeedbackFruits was a useful tool that helped me stay engaged with the material and ensured I completed all the required tasks. However, I felt that the amount of reading each week was quite extensive. I would prefer if the weekly reading load were reduced, focusing on a smaller number of papers, so that I could engage more deeply with each one.
  \item I think that I only have something to say about two of the papers that we had to read, and those would be the machine learning paper, which I found too specific and required a bunch of notions that I didn't have, and the paper Science, Security and the Elusive Goal of Security as a Scientific Pursuit, which I found harder to follow because of the large amount of quotes that did not really suit the tone of the paper. All in all, my improvement suggestion could be for the first article to add a ML cheatsheet of some sort to grasp the missing concept, and for the second article maybe pick a different one that goes more straight to the point.
  \item Material is good as is.
  \item I think the papers should be short and precise rather than very long as it would help to easily grasp the concepts.
  \item It was extremely exhausting and annoying to be forced to read countless papers throughout the course. I understand that these papers are relevant for the topics they're introduced in as preparatory material, but I still believe that the level of understanding we gain from having to read these papers is far too small for the cost imposed.
  \item Not sure.
  \item I think the amount of papers to read could be slightly more balanced, since for last seminar there was a lot less reading, as compared to for example previous seminar with 10 papers.
  \item More basic level information about the subject at study (scientific methods of cyber security), like lectures or course literature. I found it hard to abstract the knowledge just based on papers
  \item I think only reading papers as a too focused preparation. Maybe including more short essays or research on your own might be useful. However the papers selected were really interesting and helpful!
  \item I think you should state more clearly what you expect from the written reflections in terms of word count and how thorough/deep the reflections should be.  
  \item I actually thought the prep material was very well-structured and represented a good preparation for each of the upcoming seminars. Hence, I have no suggestions as to how it could be improved.
  \item One of the reflections did not accept a written response, only files. I think it was models part 2. This was slightly inconvenient as all the other ones worked fine. Otherwise I found the prep material to be good.
  \item The prep material can be improved by giving clearer instructions for the reflections and what to focus on in the annotations. Some topics needed more or better articles, while others had too many. It would also help if the papers were more clearly connected to the seminar discussions.
  \item Better instructions for reflection.  Fewer papers and do a deep dive on them instead.
  \item Some topics, like the machine learning papers, were a bit difficult to follow without the necessary prerequisites. Figuring out some way to make the topic more readable without having had a course on the subject beforehand would be nice.  Additionally, it might be good to clarify what one is meant to get out of a paper. Are we supposed to read the result sections, or is that superfluous to what the course is concerned with?
  \item Hmm. I don't really know. I think almost all of the papers were interesting to read in the context of methodologies of security research. I also thought it was nice to annotate and combine the reading with something practical. It stuck better like that.
  \item This might sound stupid, but I kind of struggled while writing the first reflections.... maybe a short guide on how to write one?
\end{enumerate}

\subsubsection{AI-Generated Summary:}

\begin{itemize}
\item \textbf{Main themes}
  \begin{itemize}
    \item Overall quality seen as good/engaging; several found the materials well-structured and FeedbackFruits helpful.
    \item Reading load is perceived as high; preference for fewer/shorter papers with deeper engagement and better week-to-week balance.
    \item Need for clearer guidance on reflections and annotations (scope, depth, word count, focus).
    \item Accessibility and prerequisites: requests for key-concept explanations and background (especially for machine learning).
    \item Curation and alignment: desire for better linkage between papers and seminar discussions; interest in having options/choice of papers.
    \item Minor platform inconsistency noted for one reflection submission format.
  \end{itemize}

\item \textbf{Common concerns or issues}
  \begin{itemize}
    \item Excessive/imbalanced reading load leading to superficial engagement and fatigue.
    \item Ambiguity about expectations for reflections and what to extract from papers (e.g., which sections matter).
    \item Difficulty with technically demanding papers without prior coursework (notably ML).
    \item Some papers felt overly long, quote-heavy, or misaligned with the seminar.
    \item Tasks sometimes allowed completion without full reading, reducing depth.
    \item One Canvas item only accepted file uploads, unlike others.
  \end{itemize}

\item \textbf{Suggestions for improvement}
  \begin{itemize}
    \item Reduce and rebalance readings: fewer papers per week, prioritize shorter/clearer papers, and do deep dives; maintain consistent load across seminars.
    \item Provide clearer scaffolding: explicit reflection rubric (word count, depth, examples), annotation prompts, and per-paper guidance on key takeaways and sections to prioritize.
    \item Add prerequisite support: glossaries/cheat sheets (especially for ML), brief primers or recommended textbook chapters/mini-lectures on core concepts and scientific methods in cybersecurity.
    \item Refine/expand selection: replace dense or quote-heavy papers with more direct alternatives; improve alignment to seminar topics; offer optional/alternative readings or allow student-selected relevant papers; include some short essays or small independent research tasks.
    \item Adjust FeedbackFruits prompts to encourage comprehensive engagement across a paper.
    \item Ensure consistent, convenient submission options in Canvas for all reflections.
  \end{itemize}

\item \textbf{Overall sentiment}
  \begin{itemize}
    \item \textit{Mixed-positive}: Many appreciated the relevance, structure, and annotation workflow, but a sizable group reported frustration with the volume and unevenness of readings and the lack of clear guidance. Students value the materials and tools but want a more focused, guided, and accessible prep experience.
  \end{itemize}
\end{itemize}

\subsection{564853: How can the seminars be improved?}

For example the topics discussed, the grouping into breakout rooms, shorter or 
longer times for discussions, longer seminars to cover the topics more deeply.

\subsubsection{Individual Responses (21 total):}

\begin{enumerate}
  \item I do not know this either. But I can pinpoint a think which I found really good, and that is that Daniel did bend and adapt the format of the seminars, on a seminar by seminar basis, based on the people that were attending. By this I mean that he took into account both how many we were (e.g. if we were just a few people he would not split us in breakout rooms) and the "quality" of the contributions (e.g. asking people who would not, by themselves, spontaneously talk much).
  \item The seminars generally had a good pace, but in the last one or two seminars, I felt the breakout room time was longer than necessary. But it wasn't that big of an issue. 
  \item I did the project
  \item One way to improve the seminars could be to hand of more of the lead to the students. In this way, we would have to engage more in discussions and be more active during the seminars.
  \item Sometimes we had way too long time for discussions. Maybe if the questions where a bit more detailed we could have discussed them further.
  \item I think breakout rooms provided more space to talk compared to the whole class setting, which was helpful for discussing and deepening our understanding of the topics. The amount of time allocated for discussions was appropriate, so I wouldn’t suggest making the seminars any shorter or longer.
  \item Not really much to say to be honest
  \item This is fine, maybe if we have some more interesting discussions will be good.
  \item I personally far prefer lectures, where we are taught about a given topic in a conventional way, and didn't derive much value at all from grouping up with my colleagues and discussing course topics with them.
  \item Shorter discussions in break-out rooms.
  \item I liked the group discussions on the seminar, and I really appreciate accessibility of the course material remotely and big freedom in terms of choosing time to do the assignments.
  \item I think that for me I would prefer lectures - but that is a personal preference. It was often hard for us to reach conclusions on our own
  \item I really enjoyed the discussions in the seminars, I think the one where we had to pick our break-out room for a specific topic worked the best, because each person picked the topic they were interested in which improved the quality of the discussion
  \item Even though it might be controversial, I think you should enforce stricter criteria of what it means to "participate" in the seminars. Many times, it felt like some people didn't contribute anything to the discussion, leading to only a few number of people who actively participated the the discourse. 
  \item Again, I really enjoyed taking part in the seminars and thought they constitute a great way to debate on the topics which we had read about in the preparatory material. No improvements necessary as far as I am concerned.
  \item I think the seminars worked well.
  \item The seminars could be improved by having more focused topics and clearer goals for each discussion. Sometimes the breakout room groups felt too small or unbalanced. A bit more time for discussion would help explore the topics more deeply.
  \item Post the discussion points on canvas at the start of the seminar, so we don't miss any points while in the group rooms.
  \item I found them generally very nicely formatted. Occasionally, the alloted time for a specific discussion topic was a bit too long (or too short), but that is likely difficult to correct for since each group may have a different complaint.
  \item I think the mix we had was quite good. Maybe the break-out rooms went on for a bit too long sometimes, meaning the discussion had already ended for us.
  \item I had the feeling that some of the questions were too open. Having a group of more concise questions related to the papers/topic of the seminar would be positive in my opinion
\end{enumerate}

\subsubsection{AI-Generated Summary:}

\textit{Respondents generally viewed the seminars positively, valuing adaptable facilitation and breakout discussions, while asking for clearer, more focused prompts, better-balanced groups, and tighter timing—alongside a minority preference for more lecture-style content.}

\textbf{Main themes}
\begin{itemize}
\item \textbf{Breakout rooms:} Widely seen as useful for deeper discussion, but timing and group composition mattered.
\item \textbf{Discussion design:} Desire for clearer goals and more focused, concrete questions to guide debate and help reach conclusions.
\item \textbf{Facilitation and participation:} Appreciation for adaptive facilitation; concerns about uneven participation.
\item \textbf{Format preferences:} Split views—most value discussions; a minority prefer conventional lectures.
\item \textbf{Logistics/resources:} Posting discussion points early (e.g., on Canvas) and maintaining remote accessibility were appreciated.
\end{itemize}

\textbf{Common concerns or issues}
\begin{itemize}
\item \textbf{Overlong breakout sessions:} Discussions sometimes ended early, leaving idle time.
\item \textbf{Overly open-ended prompts:} Hard to reach conclusions without clearer scope or objectives.
\item \textbf{Uneven participation:} Some members contributed little, placing the burden on a few.
\item \textbf{Group sizing/balance:} Groups occasionally too small or unbalanced; variability across groups.
\item \textbf{Value perception:} A minority felt peer discussions added limited value compared to lectures.
\end{itemize}

\textbf{Suggestions for improvement}
\begin{itemize}
\item \textbf{Timing:} Shorten or dynamically adjust breakout durations; allow early wrap-up/return when groups finish.
\item \textbf{Prompts/goals:} Use more focused, detailed questions tied to readings; state clear discussion goals and expected takeaways; post prompts on Canvas at the start (or in advance).
\item \textbf{Grouping:} Enable interest-based self-selection with instructor balancing; avoid splitting when attendance is low; aim for optimal group sizes.
\item \textbf{Participation:} Set and enforce participation criteria; assign rotating roles (facilitator, scribe, spokesperson); continue inviting quieter voices.
\item \textbf{Format blend:} Add brief mini-lectures or whole-class syntheses to frame topics and consolidate conclusions; consider occasional lecture-oriented sessions.
\end{itemize}

\textbf{Overall sentiment}
\begin{itemize}
\item \textbf{Mostly positive:} Many found the seminars well-paced and effective; several reported no changes needed.
\item \textbf{Targeted tweaks:} Most requests focus on tightening breakout timing, clarifying prompts/goals, and improving participation balance.
\item \textbf{Divergent preferences:} A minority prefer traditional lectures, while others strongly value interactive discussion.
\end{itemize}

\subsection{564854: What was the best part of the seminars and prep?}

For example a favourite paper, an eye-opening topic, or just to discuss how 
this relates to TAMOS.

\subsubsection{Individual Responses (22 total):}

\begin{enumerate}
  \item I really enjoyed the paper that digressed about the history of the scientific method! It is very important to understand that even our notion of "what is knowledge" is dynamic, and evolves as time passes.
  \item I liked the paper "Tor: The Second-Generation Onion Router" because I learned some concepts that I hadn't come across before. It was interesting to view all the papers through the lens of whether their results would have external validity based on the methodology.
  \item I did the project
  \item I enjoyed reading the paper on Machine learning in Cybersecurity research. It was a paper that I felt was extra relevant with current research trends which is something I appreciate while reading papers.
  \item I had never heard about graphical passwords (except the old android pattern passwords) and I really liked the paper about them.
  \item The papers about password management was really eye-opening. It showed me that passwords don’t need to have super strict rules or be changed all the time. That was a cool insight and made me think differently about how I would approach that in the future.
  \item I really enjoyed writing the last reflections, creating our own hypothesis and applying the TAMOS concepts to a scenario that honestly I felt more interesting than the ones discussed in TAMOS.
  \item The papers about Tor were good.
  \item I did not enjoy either the seminars or the prep.
  \item Some papers were more difficult to read than others. Those that were easy was best but I don't really remeber which those are...
  \item My favorite assignment was the last one, since it allowed for most freedom in choosing topics to analyze.
  \item It was interesting to see examples of cyber security research
  \item Seeing the more practical side of TAMOS was really helpful, but also just understanding how those topics of TAMOS relate to Cybersecurity, and it's specifics. Having the field which is still evolving a lot and challenging the concepts from science is really interesting!
  \item The best part was that we got to, and learn how to, apply what we learned in TAMOS in the context of cyber security. It shed some light on practically the entirety of the TAMOS course. It previously felt very hard to draw parallells to cyber security from several parts of the TAMOS content on its own.
  \item I thought the last seminar where we prepared and discussed methodologies for different security related research topics was really interesting and good practice for our masters thesis
  \item I definitely enjoyed reading some of the papers related to encryption in the preparatory materials, such as "Why Johnny Can't Encrypt". Moreover, reading about interactions of non tech-savvy people with cybersecurity tools, constructs and primitives reinforced my opinion that more should be done to ensure that the average person can deal with the basics of cybersecurity and improve their tech-literacy. 
  \item I enjoyed discussing how TAMOS related to security as I found such discussions to be one of the main things lacking in TAMOS.
  \item The best part was discussing how the topics connect to TAMOS and real-world security issues. It was interesting to hear different perspectives in the breakout rooms, and some papers gave new insights, especially on human factors in security.
  \item Small group discussions, hearing how other people analyze the topics.
  \item I really liked the final paper since it put much of the course into perspective. 
  \item I enjoyed the module about passwords and how to conduct research about the security of passwords. It opened my eyes to how difficult it can be to do proper research and evaluation of a topic that seems easy at first.
  \item Probably what I liked the most was the part in the course where we are thaught that a system can built upon provably secure elements within a formal model yet practically insecure due to overlooked holistic considerations and/or flawed assumptions. I particularly remember the paper "On the Security of ElGamal Based Encryption" to be quite enlightening about this.
\end{enumerate}

\subsubsection{AI-Generated Summary:}

\textbf{Main themes}
\begin{itemize}
\item \textit{Applying TAMOS to cybersecurity}: Strong appreciation for connecting TAMOS concepts to concrete security problems; helped clarify previously abstract material.
\item \textit{Methodology and research rigor}: Focus on external validity, study design, and evaluation challenges; last seminar seen as good practice for thesis work.
\item \textit{Human factors and usable security}: Interest in passwords/password management, graphical passwords, and usability studies (e.g., \textit{Why Johnny Can’t Encrypt}); recognition that non–tech-savvy users need better support.
\item \textit{Cryptography and the theory–practice gap}: Insights that formally secure components can yield insecure systems due to flawed assumptions (e.g., \textit{On the Security of ElGamal Based Encryption}).
\item \textit{Network privacy/anonymity}: Positive reception of Tor papers (e.g., \textit{Tor: The Second-Generation Onion Router}).
\item \textit{Current trends and relevance}: Engagement with machine learning in cybersecurity and contemporary research examples.
\item \textit{Discussions and autonomy}: Value in small-group discussions, hearing diverse perspectives, and open-ended final assignments/reflections with freedom to choose topics and craft hypotheses.
\end{itemize}

\textbf{Common concerns or issues}
\begin{itemize}
\item Variability in reading difficulty; some papers were hard to read without guidance.
\item At least one respondent did not enjoy seminars or prep; indicates uneven engagement.
\item Perceived lack of explicit security linkage in the core TAMOS course (this module helped but highlights a gap).
\item Some found course-provided scenarios less interesting than those they devised themselves.
\end{itemize}

\textbf{Suggestions for improvement}
\begin{itemize}
\item Strengthen explicit, ongoing links between TAMOS concepts and cybersecurity via applied cases and real-world studies.
\item Maintain/expand small-group discussions and peer analysis in seminars.
\item Keep and potentially broaden open-ended, student-choice assignments that involve hypothesis formation and methodological application.
\item Curate more accessible readings; provide reading guides, summaries, or focus questions emphasizing methods and threats to validity.
\item Increase coverage of usable security and human factors (password policies, graphical passwords, encryption usability) alongside technical topics like Tor and ML.
\item Add methodology workshops/tools (e.g., templates for external validity assessment, evaluation design, and holistic threat modeling bridging formal models to practical assumptions).
\item Offer optional tiers or supplementary resources to accommodate different desired depths/reading difficulty.
\end{itemize}

\textbf{Overall sentiment}
\begin{itemize}
\item Predominantly positive. Respondents valued the practical application of TAMOS, methodological training, and discussions that offered new insights (especially on human factors). A minority reported low enjoyment or difficulty with readings.
\end{itemize}

\subsection{564856: How can the project assignment be improved?}

\subsubsection{Individual Responses (20 total):}

\begin{enumerate}
  \item Given I did not have to do the project, I cannot contribute much here.
  \item It is already good
  \item I don't think there is much to improve with the project assignment.
  \item Did not do it, don't know
  \item I have not done the project assignment :)
  \item .
  \item Perhaps clearer instructions because there was so much research on Signal that it was not feasible to cover everything. Issues with feedback fruits hand-in was a problem.
  \item I didn't do the project assignment (I don't know if this question is mandatory)
  \item N/A
  \item n/a
  \item -
  \item I did not do the project, since I participated in the seminars and did all the preparation assignments
  \item -
  \item I did not do the project assignment as my aim has always been to pass by participation.
  \item Don't know as I didn't do it.
  \item I didn’t do the project since I participated during the course but from what I saw it could be improved by giving more examples and clearer instructions especially on what is expected in the final result
  \item Did not do project
  \item N/A
  \item I don't know, I did not do the project assignment.
  \item I did not have to do it
\end{enumerate}

\subsubsection{AI-Generated Summary:}

\textbf{Summary (n=20)}

\begin{itemize}
\item \textbf{Main themes}
  \begin{itemize}
  \item Most respondents did not complete the project assignment (approximately 80\%), often because they passed via participation instead.
  \item Limited but mixed feedback from those who commented: some say it is already good; others ask for clearer guidance and examples.
  \item Noted technical issues with the submission platform (FeedbackFruits).
  \end{itemize}

\item \textbf{Common concerns or issues}
  \begin{itemize}
  \item \textit{Clarity}: Unclear instructions and expectations for the final deliverable.
  \item \textit{Scope}: Topic felt too broad (too much research to cover, e.g., on Signal), making it hard to be comprehensive.
  \item \textit{Logistics}: Problems with FeedbackFruits hand-in.
  \item \textit{Process transparency}: Uncertainty about whether the project and/or this question were mandatory.
  \end{itemize}

\item \textbf{Suggestions for improvement}
  \begin{itemize}
  \item Provide clearer, more detailed instructions, including explicit criteria for the final result and, if possible, a brief rubric.
  \item Share exemplar projects or templates to illustrate expected scope and depth.
  \item Narrow or better bound the topic scope (or offer focused prompts) to keep workload feasible.
  \item Address submission workflow: test FeedbackFruits, offer guidance, or provide an alternative submission channel.
  \item Communicate clearly the alternatives for passing (participation vs. project) and the optionality of related survey questions.
  \end{itemize}

\item \textbf{Overall sentiment}
  \begin{itemize}
  \item Predominantly neutral due to non-participation; among the few who provided content, sentiment is mildly positive (some say no improvements needed) with actionable requests around clarity, scope, and submission logistics.
  \end{itemize}
\end{itemize}

\subsection{564857: Do you have anything else you would like to add to the evaluation of the course?}

\subsubsection{Individual Responses (19 total):}

\begin{enumerate}
  \item Thank you for a nice course. I was terrified at the beginning by the seemingly convoluted iter to get a pass, but in the end it was totally worth it.
  \item I think it's better to participate in the seminars rather than just do the project. I did both, and I found it more rewarding to read the papers and then discuss them in the seminar. As the course progresses, the seminars become more engaging, and you're able to participate more actively. It's also interesting to hear perspectives you hadn't really considered while reading the papers.
  \item No
  \item .
  \item I think there should be one compensation seminar in the course if you miss one of the seminars. It is unfair to have to re do the whole course if you miss one seminar. Also, there are other courses with mandatory activities, and I had to chose to skip one of them to participate in one of the seminars. Luckily the other course did have a compensation assignment, but I had to do more work.
  \item .
  \item .
  \item Good course and good material! While it is frustrating to miss a seminar or two, I appreciate that there is another way to pass the course through the project.
  \item No
  \item no
  \item I really enjoy the course, having the seminars online makes it easy to combine with other courses. The concept and the topics studied is really eye-opening and fun to dive more into the science part of Cybersecurity. Thank you for the course!
  \item -
  \item I would like to formally thank everybody involved in organizing the course for a very educating experience that allowed me to learn a lot more about scientific papers and the methodologies employed in said papers.
  \item From what I heard from other taking the course during our quarterly cybsam seminar some people seem to have been under the impression that you were supposed to closely read every paper in its entirety, leading to them complaining about the workload of the course. Maybe further clarifying how you are supposed to engage with the material would be a good idea to avoid similar sentiments in the future. (Or maybe I misunderstood and we were supposed to closely read every paper in its entirety.)
  \item The course was interesting and I liked the connection to real-world security topics. Some instructions (especially for reflections and annotations) could be clearer but overall it was a good learning experience
  \item No
  \item I was sceptical at first but I came to enjoy it.
  \item No. Thank you for this course :)
  \item Nope
\end{enumerate}

\subsubsection{AI-Generated Summary:}

\textbf{Main themes}
\begin{itemize}
\item Strong overall appreciation: many found the course \textit{engaging, eye-opening, and worthwhile}, with gratitude for the learning experience.
\item Seminars add high value: reading and discussing papers increased engagement over time; online seminars improved accessibility and scheduling.
\item Alternative pathway: the project route is appreciated as a fallback/alternative to seminars.
\item Skill development: improved understanding of scientific papers and research methodologies; good connection to real-world cybersecurity topics.
\item Initial skepticism/fear shifted to positive experience as the course progressed.
\end{itemize}

\textbf{Common concerns or issues}
\begin{itemize}
\item Attendance policy perceived as strict: missing a single seminar can jeopardize passing; conflicts with other courses’ mandatory activities.
\item Frustration when a seminar is missed, despite the existence of the project route.
\item Ambiguity in expectations: uncertainty about whether to closely read every paper in full; unclear instructions for reflections and annotations.
\end{itemize}

\textbf{Suggestions for improvement}
\begin{itemize}
\item Provide a \textit{compensation} option (make-up seminar or assignment) for a missed seminar.
\item Clarify how deeply students should read/engage with each paper and provide crisper guidance for reflections and annotations.
\item Communicate pathways and expectations more explicitly (seminar participation vs. project route), including workload framing.
\item Maintain the online seminar format for flexibility.
\end{itemize}

\textbf{Overall sentiment}
\begin{itemize}
\item Predominantly positive and grateful; students found the course rewarding, increasingly engaging, and relevant.
\item Minor but notable concerns focus on attendance rigidity and clarity of instructions.
\end{itemize}


\chapter{Differences between the two rounds}

\section*{Course Summaries and Comparison}

\subsection*{vetcyb24p2 summary}

\begin{itemize}
  \item Workload and timing: Many perceived the reading load as high and materials were often released late, making planning difficult. Time-on-task varied, with a substantial share at 6–8 h/week, but notable fractions at both lower and higher bands.
  \item Prep material: Frequent concerns about too many/too long/too technical papers (56.8\% said too many). FeedbackFruits (FBF) was seen as clunky and inconsistently configured. Students requested clearer guidance, exemplars, rubrics, and earlier releases; several wanted materials focused more on methodology than technical details.
  \item Seminars: Generally positive for understanding and discussion (67.5\% and 57.5\% positive, respectively). Issues included uneven preparedness, group size/balance, occasional language mismatches, and insufficient time for complex topics. Requests for clearer prompts and occasional mini-lectures.
  \item Project: Significant confusion about scope, expectations, deliverables, and timelines; late release and holiday deadlines compounded issues. Mixed views on structure; strong calls for precise instructions, exemplars, and unified submission flow.
  \item Assessment clarity: Uncertainty about whether seminar participation alone sufficed to pass; “all-or-nothing” seminar policy discouraged continued attendance after a miss.
  \item Overall: Mixed sentiment. Students valued discussion and the TAMOS-to-security bridge, but logistics (late uploads, unclear instructions, platform friction) and workload balance were key pain points. Top suggestions: earlier releases, fewer but better-targeted papers, standardized guidance/templates, clearer pass criteria, and tuned seminar facilitation.
\end{itemize}

\subsection*{vetcyb25p4 summary}

\begin{itemize}
  \item Workload and timing: Still perceived as substantial, but more acceptable; many at 6–8 h/week (39.1\%). Some weeks felt imbalanced; calls to reduce and rebalance readings and favor shorter papers.
  \item Prep material: Overall seen as good/engaging; seeing others’ annotations was helpful (65.2\%). Reading was good prep (73.9\%); reflections also helpful (56.5\%), though instructions needed clearer scope and depth. Requests for primers/cheat sheets for technical topics (e.g., ML), clearer per-paper focus, and better alignment to seminar discussions.
  \item Seminars: Strongly positive for discussion (78.3\%) and understanding (65.2\%). Breakout rooms worked well but sometimes ran too long; prompts could be more focused. A minority preferred lecture-style delivery; others liked adaptive facilitation and topic-based self-selection.
  \item Project: Most did not need it (87.0\%). Limited feedback suggested clearer instructions, examples, and scope bounding; a few platform submission issues.
  \item Assessment clarity: Participation route was commonly used and appreciated; some requested a compensation option for a missed seminar and clearer depth expectations for readings.
  \item Overall: Mixed-positive to positive. Students highlighted valuable application of TAMOS to security, useful discussions, and relevant papers (human factors, Tor, ElGamal, ML). Targeted improvements: trim and balance readings, sharpen reflection/annotation guidance, and fine-tune seminar prompts and timing.
\end{itemize}

\subsection*{Comparison}

\begin{itemize}
  \item Improvements from vetcyb24p2 to vetcyb25p4:
    \begin{itemize}
      \item Prep perception: Fewer reported “too many papers” (56.8\% \textrightarrow{} 39.1\%); reflections were seen as more helpful (41.0\% \textrightarrow{} 56.5\%).
      \item Seminars: Stronger endorsement of peer discussion (57.5\% \textrightarrow{} 78.3\%) and fewer complaints about unprepared peers (22.5\% \textrightarrow{} 13.0\%). Adaptive facilitation and topic-based grouping were appreciated.
      \item Assessment pathway: Clearer uptake of the participation route (majority avoided the project), with fewer concerns about pass-by-participation ambiguity.
    \end{itemize}
  \item Persisting issues in both runs:
    \begin{itemize}
      \item Reading load and balance: Desire for fewer/shorter papers and more consistent weekly workload.
      \item Guidance: Need for clearer reflection and annotation expectations, per-paper focus, and stronger links between prep and seminar use.
      \item Seminar design: Occasional overlong breakout timings; preference for more concrete prompts and explicit takeaways; minority preference for more lecture input.
      \item Platform/logistics: Some ongoing confusion or friction with FBF and submission flows; calls for all instructions centralized and timely.
    \end{itemize}
  \item Net effect:
    \begin{itemize}
      \item vetcyb25p4 shows a notable shift toward higher satisfaction with seminars and prep utility, fewer reports of overload, and smoother assessment via participation.
      \item The remaining gaps are more about calibration (load, clarity, timing) than fundamental design, suggesting incremental refinements can further improve the experience.
    \end{itemize}
\end{itemize}

\end{document}
